\documentclass[UTF8,a4paper,12pt]{ctexart}
\usepackage{geometry,amsmath,booktabs,hyperref}
\geometry{left=2.5cm,right=2.5cm,top=2.5cm,bottom=2.5cm}
\title{高性能芯片热管理系统问题二：代理模型}
\author{数学建模项目组}\date{\today}
\begin{document}\maketitle

\section{问题与数据}
附件2包含84组样本。每组样本以针肋宽度比 $x_1$、歧管深高比 $x_2$、针肋排数 $x_3$ 为输入，以无量纲热阻 $R$、无量纲压降 $P$ 和温度非均匀性 $U$ 为输出。为刻画结构参数到性能指标的映射，分别建立三个二次响应面代理模型。

\section{二次响应面模型}
对任意性能指标 $y\in\{R,P,U\}$，取
\begin{equation}
\hat y=\beta_0+\beta_1x_1+\beta_2x_2+\beta_3x_3+\beta_{11}x_1^2+\beta_{22}x_2^2+\beta_{33}x_3^2+\beta_{12}x_1x_2+\beta_{13}x_1x_3+\beta_{23}x_2x_3.
\end{equation}
记特征向量
\[
\boldsymbol\phi=(1,x_1,x_2,x_3,x_1^2,x_2^2,x_3^2,x_1x_2,x_1x_3,x_2x_3)^\mathrm T,
\]
则 $\hat y=\boldsymbol\phi^\mathrm T\boldsymbol\beta$。采用最小二乘法
\begin{equation}
\hat{\boldsymbol\beta}=(\Phi^\mathrm T\Phi)^{-1}\Phi^\mathrm Ty.
\end{equation}

\section{拟合结果}
由脚本 ``问题二代理模型.py'' 计算得到模型系数如下。系数顺序均为 $(1,x_1,x_2,x_3,x_1^2,x_2^2,x_3^2,x_1x_2,x_1x_3,x_2x_3)$。

\begin{table}[h]
\centering\caption{三个二次代理模型的系数}
\small
\begin{tabular}{c|rrr}
\toprule
项 & $R$ & $P$ & $U$\\\midrule
$\beta_0$&0.6088059&0.3905371&1.2572195\\
$\beta_1$&-0.0461039&-0.2497283&-0.8342334\\
$\beta_2$&0.0677676&-0.1248854&-0.1931900\\
$\beta_3$&0.0004886&0.0018008&-0.0060769\\
$\beta_{11}$&0.5198733&0.6065597&1.4492512\\
$\beta_{22}$&-0.0069104&0.0126592&0.0223274\\
$\beta_{33}$&0.0001876&-0.0003405&0.0003564\\
$\beta_{12}$&-0.0218153&-0.0011769&0.0625118\\
$\beta_{13}$&-0.0187149&0.0361245&0.0052244\\
$\beta_{23}$&-0.0002234&-0.0000391&0.0012317\\
\bottomrule
\end{tabular}
\end{table}

因此三个映射关系分别为 $\hat R=\boldsymbol\phi^\mathrm T\hat{\boldsymbol\beta}_R$、$\hat P=\boldsymbol\phi^\mathrm T\hat{\boldsymbol\beta}_P$ 和 $\hat U=\boldsymbol\phi^\mathrm T\hat{\boldsymbol\beta}_U$，代入上表即可计算任意设计点的预测值。

\section{拟合与预测效果评价}
采用决定系数、均方根误差和平均绝对误差：
\begin{align}
R^2&=1-\frac{\sum_i(y_i-\hat y_i)^2}{\sum_i(y_i-\bar y)^2},\\
RMSE&=\sqrt{\frac1n\sum_i(y_i-\hat y_i)^2},\\
MAE&=\frac1n\sum_i|y_i-\hat y_i|.
\end{align}
训练集拟合与固定随机种子下的5折交叉验证结果如下。

\begin{table}[h]
\centering\caption{代理模型精度}
\begin{tabular}{c|ccc|ccc}
\toprule
&\multicolumn{3}{c|}{训练集}&\multicolumn{3}{c}{5折交叉验证}\\
指标&$R^2$&RMSE&MAE&$R^2$&RMSE&MAE\\\midrule
热阻 $R$&0.7965&0.00505&0.00421&0.7180&0.00595&0.00491\\
压降 $P$&0.9763&0.00425&0.00357&0.9674&0.00498&0.00419\\
温度非均匀性 $U$&0.7968&0.00871&0.00708&0.7460&0.00974&0.00790\\
\bottomrule
\end{tabular}
\end{table}

\section{结果分析}
压降模型的交叉验证 $R^2=0.9674$，说明二次响应面能够很好地描述压降随结构参数的变化。热阻和温度非均匀性模型的交叉验证 $R^2$ 分别为0.7180和0.7460，预测误差仍较小，但说明二者存在更强的非线性或参数交互作用。训练集与交叉验证指标差距有限，未出现明显过拟合。

因此，本模型适合用于问题三的快速优化和问题四的权重/鲁棒性分析；对预测点应限制在附件2样本的参数范围内，避免对二次响应面进行远距离外推。

\section{结论}
以三个输入参数的二次项和两两交互项建立的响应面模型，给出了结构参数到三个性能指标的显式映射。模型对压降的拟合和预测最准确，对热阻和温度非均匀性的预测达到可用于优化的精度。后续优化时采用三个代理模型联合评价，并通过交叉验证误差给出预测可信度。

\end{document}
